% ----------------------
\subsection{Section 122 (Circa 22 March 1839)}
% ----------------------
	\label{DC122}
	
	% -----
	\subsubsection{The original source, and table of verse locations in the original source}
	% -----
	
		(Note that section 122, in the original letter, immediately follows the end of section 121: ``without compulsory means it shall flow [un]to thee for ever and ever. The ends of the Earth shall [enq]uire after thy na[me]...)
	
		This section of the DC is a passage drawn from a letter written in late March 1839, from a jail in Missouri. Orson Pratt selected the passages from the letter to be added to the DC in 1876, but he didn't choose everything. The following table shows the verses and then the passages in the original letters, based on my divisions in the margin.
		
		\begin{table}[H]
			\begin{tabular}{p{4cm}p{10cm}}
				\textbf{Verses in DC 122} & \textbf{Original source} \\
				1--9 & \hyperref[EMS-1839-JS-CIR-22-MAR-1839]{Letter to Edward Partridge and the Church, circa 22 March 1839} \\
				\textbf{Verses} & \textbf{Locations in original letter} \\
				1--9 & \hyperlink{EMS-1839-JS-CIR-22-MAR-1839-h}{h--l} \\
			\end{tabular}
		\end{table}
		
	% -----
	\subsubsection{Historical Background}
	% -----
		\label{DC122-HB}
		
		% --
		\paragraph{JSP}
		% --
		
			``Around 22 March 1839, JS composed this epistle in the Clay County jail in Liberty, Missouri. Addressed to Bishop Edward Partridge and the church, the letter offered counsel on a range of issues. The epistle was the second letter JS wrote to the church in March 1839, and he evidently planned to write additional general epistles that month, offering guidance and comfort to the church after the devastating setbacks of 1838 and amid church members' forced removal from Missouri in early 1839. The first epistle, addressed to the church in general and Partridge in particular, was drafted on 20 March 1839 and included meditations on persecution and the church's divine destiny. At the conclusion of the missive, JS stated that the prisoners' intent was to `continue to offer further reflections in our next epistle.' This undated letter appears to be the promised sequel. JS was the primary author of the epistle, although the other prisoners may have assisted in its composition. Like the 20 March 1839 epistle, the undated letter shifts between three rhetorical perspectives: the first-person plural of all the prisoners, the first-person singular of JS, and the voice of Deity directed to JS. Each prisoner signed the letter.
			
			``Dating the letter presents significant challenges. On 21 March 1839, JS informed Emma Smith, `I have sent an Epistle to the church,' presumably referring to the 20 March letter. He then told her that he intended to `send an other as soon as posible,' likely referring to this undated epistle. The undated letter could have been written anytime between 21 March and 6 April, the day the prisoners left Liberty for Gallatin, Missouri, to appear before the Daviess County Circuit Court. However, JS likely wrote it around 22 March, the day he composed a letter to land speculator Isaac Galland, who had offered to sell the church land in Illinois and Iowa Territory. In the 20 March 1839 epistle, JS deferred to the judgment of church leaders in Quincy, Illinois, regarding whether to accept Galland's offer. JS's thinking evidently changed by the time he wrote the letter to Galland on 22 March. In that letter, JS remarked that `the church would be wise in making the contract' and requested that Galland reserve the land for the Saints. JS used similar language in the undated epistle, stating that `the church would do well to secure to themselves' Galland's land offer.
			
			``In the epistle, JS revisited other major themes of the 20 March 1839 letter and also included new insights. In the earlier epistle, JS reflected on lessons learned from past mistakes; in the undated missive, JS further contemplated previous errors and suggested ways to avert similar problems in the future. The second epistle also contained an extended meditation on the righteous use of priesthood power; during the meditation, the perspective transitioned from the combined voice of JS and his companions addressing a general church audience to the voice of Deity addressing JS with regard to his future influence. After reviewing JS's recent arrest, forced separation from his family, and incarceration, the divine voice assured JS that his suffering would provide necessary experience. The letter then returned to the voice of JS and his companions, alternating between first-person singular and first-person plural. The epistle also instructed the Saints to prepare affidavits describing their losses in Missouri, to be submitted to government officials. The letter then concluded with an extended affirmation of the inspired nature of the United States Constitution and the principle of religious liberty.
			
			``JS presumably dictated the rough draft of the epistle, which is in the handwriting of Alexander McRae, who acted as scribe for other lengthy documents produced in the jail. This draft contains corrections by JS. When it was completed, each of the prisoners signed the epistle before it was folded in letter style and addressed to JS's wife Emma Smith in Quincy. For reasons that remain unclear, McRae then produced a fair copy that incorporated JS's revisions to the rough draft. After McRae and JS made additional minor corrections to the fair copy, the prisoners signed it, and then it was folded and addressed to Emma Smith. The fair copy is featured here because it appears to be the most complete version of the letter sent to Illinois.
			
			``The undated general epistle may have been among the `package of letters' that church member Alanson Ripley obtained from the jail on 22 March 1839. Alternatively, if the letter was not yet completed at the time of Ripley's visit, other visitors to the jail in late March and early April, including Heber C. Kimball and Theodore Turley, may have been entrusted with the letter. The letter was presumably delivered to Emma Smith, as it was addressed to her. On 11 April 1839, Hyrum Smith's wife, Mary Fielding Smith, indicated that both the 20 March 1839 epistle and the undated letter had arrived in Quincy. `We have seen the Epistols to the Church and read them several times,' she wrote to her husband. `They seem like food for the hungrey we have taken great pleasure on perusing them.' The undated epistle evidently was circulated widely among church members in Illinois, as indicated by the early copies that Edward Partridge and Albert Perry Rockwood made. An edited version of the letter was published twice in 1840---in church newspapers in Nauvoo, Illinois, and in Liverpool, England---further increasing the letter's circulation.''

	% -----
	\subsubsection{Versions}
	% -----
		\label{DC122-V}
		
		See the \href{https://drive.google.com/file/d/1goAvNz8jS4R8slYfMG_db55Ty2z_vmgK/view?usp=sharing}{sections comparison} document.
	
		\begin{compactdesc}
			\item[JSP] \hyperref[EMS-1839-JS-CIR-22-MAR-1839]{Circa 22 March 1839}
			\item[BC] ---
			\item[DC 1835] ---
			\item[DC 1844] ---
			\item[TS] 1:9 (July 1840), 131--134
			\item[MS] 1:8 (December 1840), 193--197
			\item[MS] 5:5 (October 1844), 69--72
			\item[DN] 4:5 (2 February 1854), 17
			\item[MS] 17:6 (10 February 1855), 84--88
		\end{compactdesc}

	% -----
	\subsubsection{Section 122:1} % *DC122-1 
	% -----
		\label{DC122-1}

			\footnote{%
				% \label{}%
				As I note above, in the original letter, this passage follows immediately after the material at the end of section 121.
				}%
		THE ends of the earth shall inquire after thy name, and fools shall have thee in derision, and hell shall rage against thee;

		% \begin{framed}
		% \scriptsize
			% \begin{compactdesc}
				% \item[JSP] 
				% % \item[BC] 
				% % \item[]
			% \end{compactdesc}
		% \end{framed}

	% -----
	\subsubsection{Section 122:2} % *DC122-2 
	% -----
		\label{DC122-2}

		While the pure in heart, and the wise, and the noble, and the virtuous, shall seek counsel, and authority, and blessings constantly from under thy hand.

		% \begin{framed}
		% \scriptsize
			% \begin{compactdesc}
				% \item[JSP] 
				% % \item[BC] 
				% % \item[]
			% \end{compactdesc}
		% \end{framed}

	% -----
	\subsubsection{Section 122:3} % *DC122-3 
	% -----
		\label{DC122-3}

		And thy people shall never be turned against thee by the testimony of traitors.

		% \begin{framed}
		% \scriptsize
			% \begin{compactdesc}
				% \item[JSP] 
				% % \item[BC] 
				% % \item[]
			% \end{compactdesc}
		% \end{framed}

	% -----
	\subsubsection{Section 122:4} % *DC122-4 
	% -----
		\label{DC122-4}

		And although their influence shall cast thee into trouble, and into bars and walls, thou shalt be had in honor; and but for a small moment and thy voice shall be more terrible in the midst of thine enemies than the fierce lion, because of thy righteousness; and thy God shall stand by thee forever and ever.\phantom{.}

		% \begin{framed}
		% \scriptsize
			% \begin{compactdesc}
				% \item[JSP] 
				% % \item[BC] 
				% % \item[]
			% \end{compactdesc}
		% \end{framed}

	% -----
	\subsubsection{Section 122:5} % *DC122-5 
	% -----
		\label{DC122-5}

		If thou art called to pass through tribulation; if thou art in perils among false brethren; if thou art in perils among robbers; if thou art in perils by land or by sea;

		% \begin{framed}
		% \scriptsize
			% \begin{compactdesc}
				% \item[JSP] 
				% % \item[BC] 
				% % \item[]
			% \end{compactdesc}
		% \end{framed}

	% -----
	\subsubsection{Section 122:6} % *DC122-6 
	% -----
		\label{DC122-6}

		If thou art accused with all manner of false accusations; if thine enemies fall upon thee; if they tear thee from the society of thy father and mother and brethren and sisters; and if with a drawn sword thine enemies tear thee from the bosom of thy wife, and of thine offspring, and thine elder son, although but six years of age, shall cling to thy garments, and shall say, My father, my father, why can't you stay with us?  O, my father, what are the men going to do with you?  and if then he shall be thrust from thee by the sword, and thou be dragged to prison, and thine enemies prowl around thee like wolves for the blood of the lamb;

		% \begin{framed}
		% \scriptsize
			% \begin{compactdesc}
				% \item[JSP] 
				% % \item[BC] 
				% % \item[]
			% \end{compactdesc}
		% \end{framed}

	% -----
	\subsubsection{Section 122:7} % *DC122-7 
	% -----
		\label{DC122-7}

		And if thou shouldst be cast into the pit, or into the hands of murderers, and the sentence of death passed upon thee; if thou be cast into the deep; if the billowing surge conspire against thee; if fierce winds become thine enemy; if the heavens gather blackness, and all the elements combine to hedge up the way; and above all, if the very jaws of hell shall gape open the mouth wide after thee, know thou, my son, that all these things shall give thee experience, and shall be for thy good.

		% \begin{framed}
		% \scriptsize
			% \begin{compactdesc}
				% \item[JSP] 
				% % \item[BC] 
				% % \item[]
			% \end{compactdesc}
		% \end{framed}

	% -----
	\subsubsection{Section 122:8} % *DC122-8 
	% -----
		\label{DC122-8}

		The Son of Man hath descended below them all.%
			\footnote{%
				% \label{}%
				See \hyperref[DC88-6]{DC 88:6}.
				}
		Art thou greater than he?

		% \begin{framed}
		% \scriptsize
			% \begin{compactdesc}
				% \item[JSP] 
				% % \item[BC] 
				% % \item[]
			% \end{compactdesc}
		% \end{framed}

	% -----
	\subsubsection{Section 122:9} % *DC122-9 
	% -----
		\label{DC122-9}

		Therefore, hold on thy way, and the priesthood shall remain with thee; for their bounds are set, they cannot pass.  Thy days are known, and thy years shall not be numbered less; therefore, fear not what man can do, for God shall be with you forever and ever.

		% \begin{framed}
		% \scriptsize
			% \begin{compactdesc}
				% \item[JSP] 
				% % \item[BC] 
				% % \item[]
			% \end{compactdesc}
		% \end{framed}